\section*{Capítulo \ref{ch:b1:angular-momentum} --- Momento angular}

\begin{solution}{exo:b1:angular-momentum:1}
$Fd = 50 \times 0.30 = \qty{15}{N.m}$; $Fd\sin 60^\circ = \qty{13}{N.m}$;
a lo largo del mango, la recta de acción pasa por la tuerca: $0$.
\end{solution}

\begin{solution}{exo:b1:angular-momentum:2}
$L = mrv = 5.97\times10^{24} \times 1.50\times10^{11} \times 2.98\times10^4 =
\qty{2.7e40}{kg.m^2/s}$; $\dd A/\dd t = L/2m = rv/2 = \qty{2.2e15}{m^2/s}$.
\end{solution}

\begin{solution}{exo:b1:angular-momentum:3}
$r = \ell\sin 30^\circ = \qty{0.50}{m}$, $\omega = \qty{3.4}{rad/s}$:
$L_z/m = r^2\omega = \qty{0.84}{m^2/s}$. La recta de acción de la tensión
pasa por el pivote (\omterm{def:b1:angular-momentum:moment}{momento} nulo respecto de cualquier eje que pase por
él); el peso es paralelo al eje vertical (\omterm{def:b1:angular-momentum:moment}{momento} nulo respecto de él):
$L_z$ se conserva.
\end{solution}

\begin{solution}{exo:b1:angular-momentum:4}
$J_1 = 4.0 + 2 \times 3.0 \times 0.64 = \qty{7.8}{kg.m^2}$, $J_2 = 4.0 +
2 \times 3.0 \times 0.04 = \qty{4.2}{kg.m^2}$; $\omega_2 = \omega_1J_1/J_2 =
1.85\omega_1$: \qty{3.7}{} vueltas por segundo. $E_k = L^2/2J$: cociente
$J_1/J_2 = 1.85$; el trabajo lo hacen los músculos, que recogen los brazos
contra la tendencia centrífuga.
\end{solution}

\begin{solution}{exo:b1:angular-momentum:5}
$L_\Delta = m\ell^2\dot\theta$; tensión: \omterm{def:b1:angular-momentum:moment}{momento} nulo (corta al eje);
peso: $-mg\ell\sin\theta$ (\omterm{def:b1:angular-momentum:moment}{brazo de palanca} $\ell\sin\theta$, recuperador).
$m\ell^2\ddot\theta = -mg\ell\sin\theta$.
\end{solution}

\begin{solution}{exo:b1:angular-momentum:6}
$r_pv_p = r_av_a$: $v_p/v_a = r_a/r_p = 1.034$; con $(v_p + v_a)/2 \approx
\qty{29.8}{km/s}$: $v_p = \qty{30.3}{km/s}$, $v_a = \qty{29.3}{km/s}$.
\end{solution}

\begin{solution}{exo:b1:angular-momentum:7}
$L = mrv$ se conserva: $v_2 = v_1r_1/r_2 = \qty{4.0}{m/s}$, $\omega_2 = v_2/r_2
= \qty{16}{rad/s}$ (desde \qty{4}{rad/s}). $E_1 = \qty{0.40}{J}$, $E_2 =
\qty{1.6}{J}$: el tirón realiza \qty{1.2}{J}. La tensión es radial, pero
ahora el desplazamiento del disco tiene una componente radial (describe
una espiral hacia dentro): el trabajo $\vect T\cdot\dd\vect\ell$ ya no es nulo.
\end{solution}

\begin{solution}{exo:b1:angular-momentum:8}
\omterm{def:b1:angular-momentum:moment}{Momentos} respecto del pivote (con las longitudes medidas desde él): niño
$30 \times 1.5 =
45$ (izquierda); peso del tablón en su centro, a $0.5$ hacia la derecha: $20 \times
0.5 = 10$ (derecha); segundo niño a $x$ hacia la derecha: $45x$.
Equilibrio: $45 = 10 + 45x$, $x = \qty{0.78}{m}$. El pivote soporta el
peso total, $95 \times 9.81 = \qty{932}{N}$.
\end{solution}

\begin{solution}{exo:b1:angular-momentum:9}
Masas: $m_1a = m_1g - T_1$, $m_2a = T_2 - m_2g$; polea: $J\dot\omega =
(T_1 - T_2)R$ con $a = R\dot\omega$. Sumando: $a = (m_1 - m_2)g/(m_1 + m_2 +
J/R^2) = 9.81/(3.0 + 1.0) = \qty{2.5}{m/s^2}$ (frente a \qty{3.3}{m/s^2});
$T_1 = m_1(g - a) = \qty{14.7}{N}$, $T_2 = m_2(g + a) = \qty{12.3}{N}$ ---
distintas, porque la polea necesita un \omterm{def:b1:angular-momentum:moment}{momento} neto.
\end{solution}

\begin{solution}{exo:b1:angular-momentum:10}
$r_1v_1 = r_2v_{\perp 2}$: $v_{\perp 2} = 40 \times 0.5/10 = \qty{2.0}{km/s}$.
La \omterm{cor:b1:angular-momentum:central}{ley de las áreas} da solo la componente perpendicular; la radial exige
la conservación de la energía (\cref{ch:b1:central-forces}).
\end{solution}

\begin{solution}{exo:b1:angular-momentum:11}
$J\omega$ constante con $J \propto R^2$: $T' = T(R'/R)^2 = 25 \times 86400
\times (10^4/7\times10^8)^2 = \qty{4.4e-4}{s}$. Rapidez ecuatorial $2\pi R'/T'
= \qty{1.4e8}{m/s}$, la mitad de la velocidad de la luz: el colapso no
puede conservar todo el \omterm{def:b1:angular-momentum:L}{momento angular} --- se desprende de buena parte
(y entra en juego la relatividad).
\end{solution}

\begin{solution}{exo:b1:angular-momentum:12}
$L_z = mr^2\omega$, y $r$ cambia: $\dd L_z/\dd t = 2mr\dot r\omega \neq 0$.
Ese ha de ser el \omterm{def:b1:angular-momentum:moment}{momento} de la reacción transversal $N$ de la varilla: $rN =
2mr\dot r\omega$, luego $N = 2m\dot r\omega$ --- la varilla empuja la cuenta
de lado mientras esta se desliza hacia fuera (el término de Coriolis de
\cref{ch:b1:non-inertial-frames}).
\end{solution}

\begin{solution}{pb:b1:angular-momentum:1}
\textbf{1.} $m\omega^2r = GMm/r^2$: $\omega = \sqrt{GM/r^3} = \sqrt{3.98\times
10^{14}/5.66\times10^{25}} = \qty{2.65e-6}{rad/s}$, periodo $2\pi/\omega =
\qty{2.37e6}{s} = \qty{27.4}{d}$.

\textbf{2.} $L_{\text{orb}} = 7.35\times10^{22} \times 1.475\times10^{17}
\times 2.65\times10^{-6} = \qty{2.9e34}{kg.m^2/s}$.

\textbf{3.} $v = \sqrt{GM/r}$, $L = mrv = m\sqrt{GMr}$; $\dd L/\dd r =
\tfrac12 m\sqrt{GM/r} = L/2r$.

\textbf{4.} $J = 0.33 \times 5.97\times10^{24} \times 4.06\times10^{13} =
\qty{8.0e37}{kg.m^2}$; $\Omega = \qty{7.29e-5}{rad/s}$; $L_{\text{rot}} =
\qty{5.8e33}{kg.m^2/s}$, la quinta parte de $L_{\text{orb}}$.

\textbf{5.} $J_m \approx 0.4 \times 7.35\times10^{22} \times 3.0\times10^{12} =
\qty{8.9e34}{kg.m^2}$, por $\omega$: \qty{2.4e29}{kg.m^2/s}, $10^{-5}$ del
orbital: despreciable.

\textbf{6.} La atracción del Sol actúa sobre el centro de masas del
sistema (sin \omterm{def:b1:angular-momentum:moment}{momento} respecto de él a primer orden); su efecto de marea
sobre el \omterm{def:b1:angular-momentum:moment}{par} es pequeño: el $L$ total se conserva con la precisión que
aquí hace falta.

\textbf{7.} El abultamiento más cercano, adelantado respecto de la línea
Tierra--Luna, es atraído por la Luna con una fuerza cuya componente
tangencial se opone a la rotación terrestre.

\textbf{8.} \omterm{def:b1:angular-momentum:moment}{Momento} negativo sobre la rotación: la Tierra frena y el día
se alarga.

\textbf{9.} El abultamiento tira de la Luna hacia delante: \omterm{def:b1:angular-momentum:moment}{momento}
positivo sobre la órbita, $L_{\text{orb}}$ crece.

\textbf{10.} $L_{\text{orb}} \propto \sqrt r$ crece, luego $r$ crece; $v \propto
1/\sqrt r$ disminuye.

\textbf{11.} El tirón hacia delante realiza un trabajo positivo, que eleva
la energía de la Luna; una órbita más alta es una órbita más lenta: la
energía va a la altura (potencial), más que la \omterm{thm:b1:work-and-energy:ke}{energía cinética} perdida.

\textbf{12.} Acoplamiento por marea de la Tierra: día igual al mes, sin
abultamiento adelantado y sin \omterm{def:b1:angular-momentum:moment}{par}.

\textbf{13.} $\Delta\Omega/\Omega = -\qty{2.3e-3}{s}/\qty{86164}{s} =
\num{-2.7e-8}$ por siglo.

\textbf{14.} $\Delta L_{\text{rot}} = L_{\text{rot}}\Delta\Omega/\Omega =
\qty{-1.55e26}{kg.m^2/s}$ por siglo.

\textbf{15.} $\Delta L_{\text{orb}} = +\qty{1.55e26}{kg.m^2/s}$ por siglo.

\textbf{16.} $\Delta r = 2r\,\Delta L/L = 2 \times 3.84\times10^8 \times 1.55
\times10^{26}/2.9\times10^{34} = \qty{4.1}{m}$ por siglo, \qty{4.1}{cm/yr}:
cerca de los \qty{3.8}{cm/yr} medidos (el alargamiento del día tiene
también una parte no debida a las mareas).

\textbf{17.} $\Delta T/T = \tfrac32\Delta r/r = 1.6\times10^{-8}$: $2.37\times
10^6 \times 1.6\times10^{-8} = \qty{0.04}{s}$ por siglo.

\textbf{18.} $\Delta E_{\text{rot}} = L_{\text{rot}}\Delta\Omega = 5.8\times10^{33}
\times 7.29\times10^{-5} \times (-2.7\times10^{-8}) = \qty{-1.1e22}{J}$ por
siglo.

\textbf{19.} $\Delta E_{\text{orb}} = GMm\Delta r/2r^2 = 2.93\times10^{37}
\times 4.1/(2 \times 1.475\times10^{17}) = \qty{+4.1e20}{J}$ por siglo.

\textbf{20.} Neto $-1.1\times10^{22} + 4\times10^{20} \approx \qty{-1.06e22}{J}$
por siglo: $3.4\times10^{12}\ \unit{W}$, disipados en forma de calor por el
rozamiento de marea en los océanos y en la Tierra sólida.

\textbf{21.} La quinta parte de la potencia de la humanidad y la
catorceava del flujo geotérmico --- un calefactor pequeño pero permanente.

\textbf{22.} $L_{\text{tot}} = 3.5\times10^{34}$. Para $r = 5.5\times10^8$:
$\omega = \sqrt{GM/r^3} = \qty{1.55e-6}{rad/s}$, $J\omega = 1.2\times10^{32}$,
$m\sqrt{GMr} = 7.35\times10^{22} \times 4.68\times10^{11} = 3.44\times10^{34}$;
la suma vale $3.45\times10^{34}$: coherente. Periodo $2\pi/\omega = \qty{4.1e6}{s}
= \qty{47}{d}$.

\textbf{23.} Las mareas solares seguirán frenando la Tierra por debajo del
ritmo orbital de la Luna, de modo que el \omterm{def:b1:angular-momentum:moment}{par} Tierra--Luna se invertiría
entonces; además, el Sol habrá evolucionado mucho antes (la escala de
tiempo es de decenas de miles de millones de años).

\textbf{24.} La Tierra levantó mareas en la Luna, que frenaron su rotación
hasta que una cara quedó mirándonos; al ser más ligera y estar cerca de un
compañero más pesado, el acoplamiento de la Luna fue mucho más rápido.

\textbf{25.} Se conserva: el \omterm{def:b1:angular-momentum:L}{momento angular} total. Se intercambia:
\omterm{def:b1:angular-momentum:L}{momento angular} de la rotación terrestre a la órbita lunar, unos
\qty{1.6e26}{SI} por siglo. Se disipa: \omterm{thm:b1:work-and-energy:em}{energía mecánica}, unos pocos
teravatios, en forma de calor de marea.
\end{solution}
